%============ abstract_ar.tex — Arabic abstract ===================
\clearpage
\thispagestyle{plain}
\begin{center}
\begin{Arabic}
{\LARGE \textbf{الملخص}}
\end{Arabic}
\end{center}
\vspace{12pt}
\begin{Arabic}
يُضعف مرض باركنسون التحكم العصبي العضلي في الكلام، فيُحدث في الصوت
تغيّرات قابلة للقياس قد تظهر مبكراً في مسار المرض. يقدّم هذا المشروع
تطويراً وتقييماً لنموذج بحثي أولي متكامل وقابل لإعادة الإنتاج، يحلّل
التسجيلات الصوتية الخام للكلام المتصل بحثاً عن الأنماط الصوتية
المرتبطة بمرض باركنسون باستخدام تعلّم الآلة التقليدي القابل للتفسير،
مع تركيز خاص على التحقّق الأمين من الأداء.

تُستخلص من كل تسجيل، بعد معالجة مسبقة موحّدة وتقطيع إلى مقاطع مدتها
عشر ثوانٍ، أربع وسبعون خاصية صوتية ذات أساس فسيولوجي --- إحصاءات
التردد الأساسي، والارتجاف، والتذبذب، ونسبة التوافقيات إلى الضجيج،
وبروز القمة الكبسترالية المُنعَّم، وإحصاءات التوقّفات، ومعاملات
الكبستروم على مقياس ميل مع ديناميكياتها --- ثم يُؤخذ متوسطها لكل
تسجيل. قورنت المصنّفات على قاعدة البيانات MDVR-KCL (سبعة وثلاثون
شخصاً) بتحقّق متقاطع مستقل تماماً عن المتحدث، وخضع اختيار النموذج
لقواعد أُعلنت قبل إجراء التجارب. حقّق نموذج الغابة العشوائية المعتمد
دقة متوازنة على مستوى الشخص قدرها \textenglish{0.822} بانحراف معياري
\textenglish{0.032}، وحساسية \textenglish{0.771}، ونوعية
\textenglish{0.873}، ومساحة تحت منحنى خصائص التشغيل للمستقبِل قدرها
\textenglish{0.864}، وكان مدى التردد الأساسي وتغيّرية الطيف في
مقدمة الخصائص المؤثرة. وأظهرت مقارنة مضبوطة أن تقسيماً خاطئاً متعمداً
على مستوى المقاطع يضخّم نتيجة الأنبوب نفسه من \textenglish{0.775} إلى
\textenglish{0.909}، بما يقيس كمّياً ضرورة التحقّق المستقل عن
المتحدث. كما أسفر التحقّق الخارجي على قاعدة بيانات إيطالية مستقلة
(واحد وستون متحدثاً)، مع تجميد النموذج وتوحيد عرض النطاق الترددي
لإبطال عامل مُربك مكتشف في معدّل العيّنات، عن مساحة تحت المنحنى قدرها
\textenglish{0.701} ودقة متوازنة قدرها \textenglish{0.629} على مستوى
المتحدث عبر اختلاف اللغة وظروف التسجيل. ويعرض النموذج الأولي المصاحب
النتائج الواقعة بين \textenglish{0.35} و\textenglish{0.65} بوصفها
غير حاسمة، ويرفق تحذيرات عن جودة التسجيل، ويصوغ كل مخرجاته بوصفها
مؤشراً بحثياً غير تشخيصي يستلزم تقييماً سريرياً متخصصاً.
\end{Arabic}
\clearpage
